% ECONOMICS_PROBLEM: {"id": "C73-10", "title": "Ordinary uniform approximate equilibria in quitting games", "jel_code": "C73", "source_id": "OP-0591"}
% !TeX program = xelatex
\documentclass[11pt,a4paper]{article}
\usepackage[margin=23mm]{geometry}
\usepackage{fontspec}
\setmainfont{Latin Modern Roman}
\usepackage{amsmath,amssymb,mathtools}
\usepackage{xcolor,enumitem,booktabs,longtable,array,xurl,hyperref}
\definecolor{muted}{HTML}{53636C}
\hypersetup{colorlinks=true,urlcolor=blue,linkcolor=blue}
\setlength{\parindent}{0pt}
\setlength{\parskip}{5pt}
\newcommand{\R}{\mathbb R}
\newcommand{\E}{\mathbb E}
\newcommand{\N}{\mathbb N}
\newcommand{\ind}{\mathbf 1}
\newcommand{\Var}{\operatorname{Var}}
\newcommand{\Cov}{\operatorname{Cov}}
\newcommand{\supp}{\operatorname{supp}}
\DeclareMathOperator*{\argmin}{arg\,min}
\DeclareMathOperator*{\argmax}{arg\,max}
\newcommand{\entrymeta}[3]{{\sffamily\small\color{muted}\textbf{\texttt{#1}}\quad|\quad #2\par #3\par}}
\newcommand{\Problemref}[1]{\texttt{#1}}
\newcommand{\JELref}[2]{\texttt{#2} (JEL \texttt{#1})}
\begin{document}
\section*{Ordinary uniform approximate equilibria in quitting games}
\textbf{Problem ID:} \texttt{C73-10}\quad \textbf{JEL:} \texttt{C73}\par
% BEGIN ECONOMICS STATEMENT
\label{problem:OP-0591}
\label{supp:quitting-games}

\noindent\textbf{Original sources:} \emph{Quitting Games and Uniform
	Equilibria} and the quitting-game components of \emph{Uniform
	Equilibrium in Finite Stochastic Games}, both dated 8 October 2026.
\textbf{Formulation:} Separate open problem, with terminal,
uniform, finite-certificate, and refinement variants.

\subsection*{Definitions and assumptions}

Fix $n\ge2$ players $I=\{1,\ldots,n\}$. At every nonabsorbing
stage, each player independently chooses \emph{continue} $C_i$ or
\emph{quit} $Q_i$. Let $\theta\in\mathbb N_0\cup\{\infty\}$ be the
first stage at which at least one player quits. If the nonempty set of
quitters at that stage is $J\subseteq I$, the game absorbs forever with
stage payoff $r(J)\in\R^n$. If no player ever quits, the stage payoff
remains $c=r(\varnothing)\in\R^n$. The entire payoff table
$\{r(J):J\subseteq I\}$, including simultaneous-quitting payoffs, is
part of the specification. No public external randomization or
correlation device is permitted.

A behavioral profile $x=(x_i)_i$ assigns conditional quitting
probabilities after each survived stage. Equivalently for this model,
one may work with independent distributions of each player's planned
quitting time $T_i\in K:=\mathbb N_0\cup\{\infty\}$; these induce the
same first-quitting outcomes. Let $J_\theta$ be the set of players
quitting at the first finite quitting stage and write
\[
\Gamma_i^\infty(x)=\E_x\!\left[
r_i(J_\theta)\mathbf1_{\{\theta<\infty\}}
+c_i\mathbf1_{\{\theta=\infty\}}\right].
\]
Define terminal strategic regret and its optimal value by
\[
\begin{aligned}
	\operatorname{Reg}_r(x)
	&=\max_{i\in I}\sup_{\widetilde x_i}
	\bigl[\Gamma_i^\infty(\widetilde x_i,x_{-i})
	-\Gamma_i^\infty(x)\bigr],\\
	g(r)&=\inf_{x}\operatorname{Reg}_r(x).
\end{aligned}
\]
For a horizon $H$ and $0<\lambda<1$, let
$\gamma_i^H(x)$ and $\gamma_i^\lambda(x)$ denote respectively
expected $H$-stage averages and normalized $\lambda$-discounted sums
of the pre-absorption stage payoff $c$ followed by $r(J_\theta)$.

For finite-grid formulations let $k\ge1$,
$K_{nk}=\{0,\ldots,nk-1,\infty\}$, and $\mathcal G_{n,k}$ be the
set of product stopping-time laws with each marginal an average of
$k$ point masses in $K_{nk}$ (repetitions permitted). Set
\[
m_k(r)=\min_{x\in\mathcal G_{n,k}}\operatorname{Reg}_r(x),\quad
D(r)=\max_{i\in I}\left(\max_{J\subseteq I}r_i(J)
-\min_{J\subseteq I}r_i(J)\right).
\]

\subsection*{Formal research question}

\begin{enumerate}[label=\textbf{(\alph*)},leftmargin=2.5em]
	\item \textbf{Ordinary terminal/undiscounted approximate Nash existence.}
	Is the universal assertion
	\[
	\boxed{\quad \forall n\ge2\ \forall r:\quad
		g(r)=0\quad}
	\]
	true? Equivalently, does every such game admit an ordinary behavioral
	$\varepsilon$-Nash equilibrium for $\Gamma^\infty$ for every
	$\varepsilon>0$? An ordinary Nash equilibrium need not be stationary
	or exact, and its strategies may depend on $\varepsilon$.
	
	\item \textbf{Uniform finite-horizon and discounted version.}
	For every $r$ and $\varepsilon>0$, does there exist a single
	behavioral profile $x^\varepsilon$ and thresholds
	$H_\varepsilon<\infty$, $\lambda_\varepsilon>0$ such that, for every
	$i$, every unilateral behavioral deviation $\widetilde x_i$,
	$H\ge H_\varepsilon$, and $0<\lambda\le\lambda_\varepsilon$,
	\[
	\begin{aligned}
		\gamma_i^H(\widetilde x_i,x_{-i}^\varepsilon)
		-\gamma_i^H(x^\varepsilon)&\le\varepsilon,\\
		\gamma_i^\lambda(\widetilde x_i,x_{-i}^\varepsilon)
		-\gamma_i^\lambda(x^\varepsilon)&\le\varepsilon?
	\end{aligned}
	\]
	The supplied quitting-game note records equivalence of universal
	existence in (a) and (b) for this model; they remain distinct
	\emph{definitions}, and no subgame-perfect conclusion is implicit.
	
	\item \textbf{Finite-grid reformulation and counterexample certificate.}
	Determine the truth of the universal exact inequality
	\[
	\boxed{\quad
		m_k(r)\le\frac{(2n-1)D(r)}{k}
		\quad\text{for all }n,r,k\ge1\quad}
	\]
	(where $n\ge2$ and $k\ge1$). For the opposite direction, find a
	complete payoff table $r$ and finite $k$ for which
	\[
	m_k(r)>\frac{(2n-1)D(r)}{k}.
	\]
	The supplied note establishes the finite-grid reduction and treats
	the universal inequality as an \emph{unproved} equivalent target.
	A rigorous counterexample must permit arbitrary behavioral
	strategies, including nonstationary and nonabsorbing play.
	
	\item \textbf{Stronger sequential refinements and restricted classes.}
	Separately ask for $\varepsilon$-subgame-perfect equilibria, with
	incentive inequalities after \emph{each} finite survival history,
	and for equilibria under stationary or bounded-memory restrictions.
	Determine whether complete payoff tables with four or more players,
	including those whose singleton-quitting matrices lie in a difficult
	complementarity-index regime, always admit ordinary equilibria.
	Any absorption-path classification must allow both discrete jumps
	(simultaneous quitting at a stage) and continuous small-hazard
	segments, rather than only a restricted one-active-player path class.
\end{enumerate}

\subsection*{Interpretation and scope}

Quitting games are a special one-nonabsorbing-state subclass of the
stochastic games in \Problemref{OP-0590}; they are retained here as a
\emph{separate} problem because the independent stopping-time,
complementarity, and finite-grid formulations are specific to them.
Existence of a sunspot or correlated approximate equilibrium does not
answer (a) or (b), which allow no extra correlating signal. Failure of
exact finite-terminal equilibria to approximate the infinite game,
nonexistence of stationary equilibria, or nonexistence in a restricted
absorption-path family is not by itself a negative answer.

\subsection*{Source audit and status}

The supplied sources report positive results for two- and three-player
quitting games and for several restricted payoff/matrix classes, while
identifying the unrestricted ordinary existence problem for four or
more players as unresolved. They also report sunspot approximate
equilibrium existence in all finite-player quitting games; this is a
different solution concept. The finite-grid equivalence and error
bound are results claimed and argued in the supplied working note,
not independently peer-reviewed here. Neither a universal proof of
the inequality in (c) nor a certified violating table is supplied.

\subsection*{References}

\begin{enumerate}
	\item[S1] E. Solan and N. Vieille (2001).
	\emph{Quitting games}. Mathematics of Operations Research 26(2),
	265--285. \url{https://doi.org/10.1287/moor.26.2.265.10549}.
	\item[S2] E. Solan and N. Vieille (2002).
	\emph{Quitting games---an example}.
	International Journal of Game Theory 31, 365--381.
	\url{https://doi.org/10.1007/s001820200125}.
	\item[S3] E. Solan and O. N. Solan (2020).
	\emph{Quitting games and linear complementarity problems}.
	Mathematics of Operations Research 45(2), 434--454.
	\url{https://doi.org/10.1287/moor.2019.0996}.
	\item[S4] G. Ashkenazi-Golan, I. Krasikov, C. Rainer, and E. Solan (2024).
	\emph{Absorption paths and equilibria in quitting games}.
	Mathematical Programming 203, 735--762.
	\url{https://doi.org/10.1007/s10107-022-01807-6}.
	\item[S5] G. Ashkenazi-Golan, I. Krasikov, C. Rainer, and E. Solan (2026).
	\emph{The APS approach for undiscounted quitting games}.
	International Journal of Game Theory 55, article 19.
	\url{https://doi.org/10.1007/s00182-026-00982-6}.
	\item[S6] Supplied research notes: \emph{Quitting Games and Uniform
		Equilibria} and the quitting-game analysis within \emph{Uniform
		Equilibrium in Finite Stochastic Games}, 8 October 2026.
\end{enumerate}

\subsection*{Common conventions: Source mathematical conventions}

\textbf{Finite games.} For a finite set $A$, $\Delta(A)$ is its probability
simplex. Players choose independent mixed strategies in Nash equilibrium;
correlation is allowed only when explicitly part of the solution concept.
In a game with payoff functions $u_i$ and strategy profile $x$, an additive
$\varepsilon$ Nash equilibrium satisfies
\[
 u_i(x)\geq u_i(x_i',x_{-i})-\varepsilon
 \quad\text{for every player $i$ and every admissible deviation $x_i'$}.
\]
For numerical approximation, payoffs are normalized as locally specified.
An $\varepsilon$ payoff residual is distinct from distance at most
$\varepsilon$ to the set of exact equilibria. Point convergence, convergence
of empirical play, and convergence in distance to an equilibrium set are
also distinct.

\textbf{Computational representations.} Unless stated otherwise, a complexity
question takes rational data with binary encoding; polynomial time is measured
in total bit length. A fixed-accuracy polynomial algorithm is not necessarily
polynomial in $1/\varepsilon$, and neither is necessarily polynomial in
$\log(1/\varepsilon)$. Succinct games and value, demand, or sampling oracles
must declare their interfaces. Exact real solutions require an explicit
representation; an unspecified list of arbitrary real numbers is not a
finite computational output. A claimed algorithmic barrier is meaningful
only for its specified model and reductions.

\textbf{Dynamic games.} Histories, information, observation, and allowable
randomization are part of the model. A stationary strategy depends only on
the current state; a positional strategy in a deterministic graph game
chooses one action at each controlled vertex. Nash equilibrium at an initial
state and equilibrium after every history are different requirements.
An expectation of a pathwise $\liminf$ payoff, a $\liminf$ of expected averages,
a limiting expected average, and a uniform finite-horizon equilibrium are
different criteria. None is silently substituted for another.

\textbf{Allocation and mechanisms.} A complete allocation assigns every item
exactly once unless otherwise stated. Zero-valued goods, negative values,
capacity constraints, feasibility, lotteries, and copies of goods affect
fairness definitions and are specified locally. Dominant-strategy incentive
compatibility, Bayesian incentive compatibility, universal truthfulness,
and truthfulness in expectation impose different inequalities. Whenever
an objective ratio has zero denominator, its convention is stated or those
instances are excluded.

\textbf{Infinite-dimensional models.} Expectations, integrals, stochastic
equations, and optimization objectives require their local regularity,
integrability, filtrations, and admissible domains. A backward--forward
mean-field system is a boundary-value problem; it is not an initial-value
semiflow without an additional construction. Long-time, large-population,
and vanishing-noise limits require an order of limits and a convergence
topology. If a minimum need not be attained, the objective is its infimum
or supremum and the approximation requirement is stated separately.
% END ECONOMICS STATEMENT
\end{document}
